\documentclass[10pt,a4paper]{amsart}
\usepackage[margin=19mm]{geometry}
\usepackage{amsmath,amssymb,mathtools}
\usepackage[scr=rsfs,cal=euler]{mathalfa}
\usepackage{xfrac}
\usepackage[unicode,hidelinks,bookmarks=false]{hyperref}
\newcommand{\ie}{i.e.\ }
\newcommand{\cf}{cf.\ }
\newcommand{\resp}{respectively\ }
\newcommand{\Subsection}{\subsection}
\providecommand{\nobreakdash}{\nobreak\hbox{-}}
\setlength{\emergencystretch}{2em}
\multlinegap0pt
\usepackage{amssymb}
\usepackage{enumerate}
\newtheorem{lemma}{Lemma}
\title{Crystallizations of generalized lens spaces}
\author{Basudeb Datta}
\date{Source: arXiv:2603.21328v1 (CC BY 4.0)}
\begin{document}
\maketitle

\noindent\emph{Source and licence.} Crystallizations of generalized lens spaces. Version arXiv:2603.21328v1 (CC BY 4.0), distributed by arXiv under the Creative Commons Attribution 4.0 International licence, http://creativecommons.org/licenses/by/4.0/, as recorded in the arXiv licence metadata for this exact version and shown on the abstract page https://arxiv.org/abs/2603.21328v1. Full licence notice: \texttt{licenses/arxiv-2603.21328-CC-BY-4.0.txt}.\par\smallskip

\noindent\textit{Editorial scope.} Counted unit: the article's lemma lem:equality (source0.tex lines 175--177). The article's complete original proof of it is delivered with it (source0.tex lines 179--205, one \texttt{proof} environment of 27 lines). No mathematics is written, repaired or re-derived: the statement and the proof are the article's own lines, in the article's own notation, and no retained line is substituted or re-flowed.  Retained uncounted as the source-local context the proof needs: lines 116--118; lines 141--144; lines 169--171.  Nothing was omitted from any retained span, so the package carries no omission record.  No retained line contains a citation, so the wrapper carries no bibliography and no bibliography entry.  No retained line contains a figure environment, an image-inclusion command or drawing code, so no image is delivered and the package declares no asset dependency.  Editorial formatting only, in addition: an AMS \texttt{amsart} preamble that restates the article's own declarations and macros used by the retained lines, namely the environments \texttt{lemma} on separate counters, together with the standard packages those lines need.  The retained environments print in this document as Lemma 1 in order of appearance, whereas the article prints the same items with the numbering of its own full document.  The complete native source of the article as distributed in the arXiv e-print for this exact version is bundled unchanged as \texttt{sources/196937/source0.tex}.
\begin{align}\label{eq:action}
h(z_0, z_1, \dots, z_n) & := (e^{2\pi i/p}z_0, e^{2\pi iq_1/p}z_1, \dots, e^{2\pi iq_n/p}z_n). 
\end{align}

\begin{align}\label{eq:Phi}
\Phi & :  |\Sigma| \to S^{\hspace{.2mm}2n+1}, \mbox{ given by } 
  w \mapsto w/\|w\|, 
 \end{align}

\begin{align}\label{eq:|rho|}
|\rho|(w) := t_1\rho(w_1)+ \cdots + t_m\rho(w_m)  
\end{align}

\begin{lemma}\label{lem:equality}
Let $h: S^{\hspace{.2mm}2n+1}  \to S^{\hspace{.2mm}2n+1} $,  $\Phi: |\Sigma| \to S^{\hspace{.2mm}2n+1}$ and $|\rho|: |\Sigma| \to |\Sigma|$ be as in \eqref{eq:action},  \eqref{eq:Phi} and \eqref{eq:|rho|} respectively. Then $\Phi\circ|\rho|\circ\Phi^{-1}\equiv h$.
\end{lemma}

\begin{proof}
Let $w\in |\Sigma|$. Suppose, $w$ is in the relative interior of $\alpha$. Let $\alpha = \alpha_1 \sqcup\cdots \sqcup \alpha_r$, where $\alpha_s$ is either a vertex or an edge of $\Sigma_{i_s}$, $s=1, \dots, r$ and $i_1, \dots, i_r$ are distinct. Then $w= a_1x_1+\cdots + a_rx_r$, where $x_1, \dots, x_r$ are in the relative interior of $\alpha_1, \dots, \alpha_r$ respectively and $a_1, \dots, a_r\in (0, 1)$ with $a_1+\cdots+a_r=1$.  

\medskip 

\noindent {\bf Claim 1.} $|\rho|(w) = a_1e^{2\pi iq_{i_1}/p} x_1+ \cdots + a_re^{2\pi iq_{i_r}/p}x_r$, where $q_0=1$. 

\medskip

Without loss, let us assume that $\alpha_1 = \{w_1, w_2\}, \dots, \alpha_{\ell} = \{w_{2\ell - 1}, w_{2\ell}\}$, $\alpha_{\ell+1}= w_{2\ell+1}, \dots, \alpha_{s} = w_{\ell+s}$. Let $x_k = b_{2k-1}w_{2k-1}+ b_{2k}w_{2k}$, where $b_{2k-1}, b_{2k}>0$ and $b_{2k-1}+ b_{2k} =1$ for $1\leq k\leq\ell$. Then $w= a_1b_1w_1 + a_1b_2w_2 + \cdots + a_{\ell}b_{2\ell}w_{2\ell}+ a_{\ell+1}w_{2\ell+1} + \cdots+ a_sw_{\ell+s}$ and $a_1b_1 + \cdots + a_{\ell}b_{2\ell}+ a_{\ell+1} + \cdots+ a_r = a_1(b_1+b_2) + \cdots + a_{\ell}(b_{2\ell-1}+b_{2\ell}) + a_{\ell+1} + \cdots+ a_r = a_1+a_2+\cdots +a_r =1$.  Then, by the definition in \eqref{eq:|rho|}, 
\begin{align*}
|\rho| (w)  &= a_1b_1\rho(w_1) + a_1b_2\rho(w_2) + \cdots + a_{\ell}b_{2\ell}\rho(w_{2\ell}) + a_{\ell+1}\rho(w_{2\ell+1}) + \cdots+ a_r\rho(w_{\ell+r}) \\
& = a_1e^{2\pi iq_{i_1}/p}(b_1w_1 + b_2w_2) + \cdots +  a_{\ell}e^{2\pi iq_{i_{\ell}}/p}(b_{2\ell-1}w_{2\ell-1} +b_{2\ell}w_{2\ell}) \\
& ~~~~ \qquad + a_{\ell+1}e^{2\pi iq_{i_{\ell+1}}/p}w_{2\ell+1}  + \cdots+ a_re^{2\pi iq_{i_{r}}/p}w_{\ell+r} \\
& = a_1e^{2\pi iq_{i_1}/p} x_1+ \cdots + a_re^{2\pi iq_{i_r}/p}x_r. 
\end{align*}
This proves Claim 1. 

\smallskip

Now, $x_s\in \mathbb{C}_{i_s}$, for $s=1,\dots, r$ and $i_1, \dots, i_r$ are distinct. Therefore, $\{a_1x_1, \dots, a_sx_r\}$ and $\{a_1e^{2\pi iq_{i_1}/p} x_1, \dots, a_se^{2\pi iq_{i_r}/p}x_r\}$ are two sets of orthogonal vectors in $\mathbb{C}^{n + 1}$. This implies that $\||\rho|(w)\|^2 = \|a_1e^{2\pi iq_{i_1}/p} x_1\|^2 + \cdots + \|a_re^{2\pi iq_{i_r}/p}x_r\|^2 = \|a_1x_1\|^2 + \cdots + \|a_rx_r\|^2=\||w\|^2$. For $1\leq s\leq r$, let $a_sx_s/\|w\|= c_se_{i_s}$, where $e_0=(1, 0, \dots, 0), e_1=(0, 1, 0, \dots, 0)$, $e_n=(0, \dots, 0, 1)$ are the standard unit vectors in $\mathbb{C}^{n+1}$. Then 
\begin{align*}
(\Phi\circ|\rho|)(w) & = \frac{a_1|\rho|(x_1) + \cdots + a_r|\rho|(x_r)}{\||q|(w)\|} = \frac{a_1e^{2\pi iq_{i_1}/p} x_1+ \cdots + a_re^{2\pi iq_{i_r}/p}x_r}{\|w\|}  \mbox{ and} \\
(h\circ\Phi)(w)& = h(c_1e_{i_1}+\cdots+c_se_{i_s}) = e^{2\pi i q_{i_1}/p}c_1e_{i_1}+\cdots+ e^{2\pi i q_{i_r}/p}c_re_{i_r}. 
\end{align*}
Therefore, $(\Phi\circ|\rho|)(w)= (h\circ\Phi)(w)$ for all $w\in |\Phi|$. Hence $\Phi\circ|\rho| \cong h\circ\Phi$. This proves the lemma. 
\end{proof}

\end{document}
